% =====================================================================
%  2. Revisão bibliográfica / Fundamentação teórica — EM ANDAMENTO
%  Explicar DevOps, observabilidade e Administração Pública com base na
%  literatura escolhida.
%  2.1 e 2.2: texto do rascunho draft/articles-selection.txt
%  (PRISMA 2020, itens 5, 8, 16a e 16b).
% =====================================================================
\section{Revisão bibliográfica}\label{sec:revisao}

% ---------------------------------------------------------------------
% PRISMA 2020, item 5: critérios de elegibilidade
% ---------------------------------------------------------------------
\subsection{Critérios de elegibilidade dos estudos}\label{sec:criterios}

Foram estabelecidos critérios de inclusão e exclusão com o objetivo de
selecionar publicações pertinentes à investigação da observabilidade
aplicada à cultura DevOps no contexto da Administração Pública.

Como critérios de inclusão, consideraram-se artigos científicos publicados
entre 2020 e 2026, nos idiomas português ou inglês, que abordassem
fundamentos, práticas, capacidades ou desafios relacionados a DevOps,
observabilidade ou transformação digital no setor público, apresentando
contribuições relevantes para os objetivos da pesquisa
(Tabela~\ref{tab:inclusao}). Os artigos devem atender a todos os critérios
pertinentes.

\begin{table}[htbp]
\centering
\caption{Critérios de inclusão}\label{tab:inclusao}
\footnotesize
\begin{tabularx}{\linewidth}{@{}lX@{}}
\toprule
\textbf{Código} & \textbf{Critério} \\
\midrule
CI1 & Artigos científicos publicados entre 2020 e 2026. \\
CI2 & Publicações em português ou inglês, com texto integral acessível. \\
CI3 & Estudos relacionados a DevOps, observabilidade ou transformação digital
      na Administração Pública. \\
CI4 & Estudos com contribuições relevantes e complementares aos eixos da
      revisão. \\
\bottomrule
\end{tabularx}
\end{table}

Foram definidos como critérios de exclusão estudos sem aderência aos eixos
temáticos, publicações restritas a soluções técnicas específicas sem
contribuição identificável para a síntese proposta, trabalhos fora do
recorte estabelecido e registros duplicados (Tabela~\ref{tab:exclusao}).

\begin{table}[htbp]
\centering
\caption{Critérios de exclusão}\label{tab:exclusao}
\footnotesize
\begin{tabularx}{\linewidth}{@{}lX@{}}
\toprule
\textbf{Código} & \textbf{Critério} \\
\midrule
CE1 & Estudos sem aderência suficiente ao problema de pesquisa. \\
CE2 & Estudos com abordagem técnica muito específica, fora do recorte
      adotado. \\
CE3 & Estudos com contribuições predominantemente sobrepostas às de outros
      trabalhos priorizados, sem acréscimo considerado essencial à síntese. \\
CE4 & Publicações fora dos critérios de período, idioma ou tipo de
      documento. \\
\bottomrule
\end{tabularx}
\end{table}

A seleção foi orientada pela pertinência dos estudos à questão de pesquisa,
considerando suas contribuições conceituais, técnicas e organizacionais.

% ---------------------------------------------------------------------
% PRISMA 2020, itens 8 (processo de seleção), 16a (seleção dos estudos)
% e 16b (estudos excluídos)
% PENDENTE (item 8): explicar como a equipe aplicou os critérios.
% PENDENTE (item 16a): fluxograma com as quantidades reais.
% ---------------------------------------------------------------------
\subsection{Avaliação dos estudos}\label{sec:avaliacao}

A Tabela~\ref{tab:aprovados} apresenta os estudos aprovados e a
Tabela~\ref{tab:reprovados}, os estudos reprovados, com a justificativa de
cada decisão.

\begin{table}[htbp]
\centering
\caption{Estudos aprovados}\label{tab:aprovados}
\footnotesize
\begin{tabularx}{\linewidth}{@{}>{\raggedright\arraybackslash}p{4.2cm}X@{}}
\toprule
\textbf{Estudo} & \textbf{Justificativa} \\
\midrule
\cite{giamattei2024}\newline \emph{Monitoring tools for DevOps and microservices}
  & Analisa ferramentas de monitoramento em ambientes DevOps e microsserviços,
    contribuindo para a compreensão das práticas de observabilidade, métricas,
    logs e traces. \\
\addlinespace
\cite{plant2025}\newline \emph{Improving DevOps team performance through
  context-capability coalignment}
  & Investiga capacidades e desempenho de equipes DevOps em organizações
    públicas, estabelecendo conexão direta entre cultura DevOps e
    Administração Pública. \\
\addlinespace
\cite{azad2023}\newline \emph{DevOps critical success factors}
  & Identifica fatores críticos de sucesso na adoção de DevOps, fundamentando
    aspectos culturais, técnicos e organizacionais relevantes para o
    estudo. \\
\addlinespace
\cite{amaro2025}\newline \emph{Mapping DevOps capabilities to the software
  life cycle}
  & Relaciona capacidades DevOps às etapas do ciclo de vida do software,
    contribuindo para a fundamentação das práticas de automação,
    monitoramento e melhoria contínua. \\
\addlinespace
\cite{figueredo2026}\newline \emph{Challenges of Digital Services in Public
  Administration in the Era of Digital Transformation}
  & Identifica barreiras institucionais, organizacionais e tecnológicas à
    transformação digital na Administração Pública, contextualizando os
    desafios de adoção de novas práticas de engenharia de software. \\
\addlinespace
\cite{amaro2026}\newline \emph{Strategic initiatives in DevOps adoption}
  & Investiga a relação entre capacidades técnicas, métricas, processos e
    iniciativas estratégicas na adoção de DevOps, contribuindo para a
    discussão sobre integração entre cultura organizacional e práticas
    tecnológicas. \\
\addlinespace
\cite{linjama2026}\newline \emph{Continuous Deployment Adoption: Insights From
  a Public Sector Implementation}
  & Apresenta um estudo de caso sobre implantação de entrega contínua no
    setor público, oferecendo evidências sobre benefícios, desafios e
    particularidades da adoção de práticas DevOps na Administração
    Pública. \\
\addlinespace
\cite{kosinska2023}\newline \emph{Toward the Observability of Cloud-Native
  Applications}
  & Apresenta um mapeamento sistemático sobre observabilidade em aplicações
    cloud-native, contribuindo para a fundamentação conceitual, abordagens
    de implementação e perspectivas de observabilidade em sistemas
    distribuídos. \\
\addlinespace
\cite{alongi2022}\newline \emph{Event-sourced, observable software
  architectures}
  & Aborda a observabilidade como propriedade arquitetural de software,
    apresentando princípios e experiências práticas que complementam a
    fundamentação técnica sobre sistemas observáveis. \\
\bottomrule
\end{tabularx}
\end{table}

\begin{table}[htbp]
\centering
\caption{Estudos reprovados}\label{tab:reprovados}
\footnotesize
\begin{tabularx}{\linewidth}{@{}>{\raggedright\arraybackslash}p{3.2cm}
                                >{\raggedright\arraybackslash}p{2.2cm}X@{}}
\toprule
\textbf{Estudo} & \textbf{Decisão} & \textbf{Justificativa} \\
\midrule
\cite{gall2022} & Não utilizar
  & Duplicidade de tratamento direto dos fundamentos conceituais de DevOps já
    escolhidos em outro artigo. \\
\addlinespace
\cite{hrusto2021} & Não utilizar
  & Investiga uma solução específica de monitoramento autônomo e filtragem de
    alertas, em vez de uma análise abrangente da adoção de
    observabilidade. \\
\addlinespace
\cite{borges2024} & Não utilizar
  & Concentra-se em experimentos e decisões técnicas de observabilidade com
    injeção de falhas, um recorte arquitetural especializado. \\
\addlinespace
\cite{gomes2025} & Texto não recuperado
  & É bastante relevante ao tema, porém não conseguimos acesso ao material
    completo. \\
\addlinespace
\cite{kumar2025} & Texto não recuperado
  & É bastante relevante ao tema, porém não conseguimos acesso ao material
    completo. \\
\bottomrule
\end{tabularx}
\end{table}

% ---------------------------------------------------------------------
% Apresentar o conceito de DevOps, a integração entre desenvolvimento
% e operações, automação, colaboração e melhoria contínua, fundamentando
% com artigos selecionados.
% ---------------------------------------------------------------------
\subsection{Cultura DevOps e seus princípios}\label{sec:devops}

% ---------------------------------------------------------------------
% Explicar o conceito de observabilidade, a utilização de métricas,
% logs e traces e sua relação com práticas DevOps.
% ---------------------------------------------------------------------
\subsection{Observabilidade de aplicações e infraestrutura}\label{sec:observabilidade}

% ---------------------------------------------------------------------
% Relacionar os conceitos anteriores com desafios do setor público,
% como sistemas legados, limitações organizacionais, necessidades de
% capacitação e continuidade dos serviços digitais, conforme as
% evidências dos artigos.
% ---------------------------------------------------------------------
\subsection{DevOps e observabilidade na Administração Pública}\label{sec:setor-publico}

% Contribuicao: fundamentacao teorica do eixo DevOps e observabilidade no setor publico.
% Incluida apos a \subsection existente no arquivo 02-revisao-bibliografica.tex.
% Nao modifica os criterios de selecao, a triagem PRISMA nem as demais secoes.
% Gomes et al. (2025): somente as informacoes verificadas no resumo editorial;
% continua classificado como "texto nao recuperado" na selecao do grupo.

A adoção de DevOps em organizações públicas depende de condições organizacionais e tecnológicas que não se limitam à automação da entrega de software. Estudos sobre equipes DevOps no setor público e sobre transformação digital governamental apontam a importância de adequar capacidades técnicas às características institucionais e às necessidades dos serviços prestados \cite{plant2025,figueredo2026}. Nessa perspectiva, a relação entre DevOps e observabilidade é examinada aqui mediante três dimensões: o alinhamento entre contexto e capacidades, as barreiras à transformação digital e o potencial diagnóstico da observabilidade em sistemas distribuídos.

\subsubsection{DevOps na Administração Pública: contexto e capacidades}

Plant et al.~\citeano{plant2025} desenvolveram e avaliaram um perfil de alinhamento entre contexto organizacional e capacidades DevOps (\emph{context-capability coalignment}) em uma equipe de tecnologia de uma organização pública dos Países Baixos. A investigação longitudinal, de 13 meses, adotou \emph{Action Design Research} e combinou questionários, oficinas, opiniões de especialistas e entrevistas. Os fatores contextuais identificados foram complexidade tecnológica relativamente elevada, forte regulamentação, interdependências departamentais e alta relevância do sistema. O perfil elaborado para aquela equipe incluiu capacidades de gestão do conhecimento, monitoramento e documentação dos sistemas, segurança e controle de acesso, gestão de requisitos, feedback dos usuários e colaboração entre equipes.

Ao final da intervenção, os pesquisadores identificaram melhora moderada no alinhamento entre contexto e capacidades e melhora percebida em aspectos do desempenho da equipe. Entretanto, as ações destinadas ao compartilhamento de logs e ao monitoramento automatizado de eventos foram adiadas, por dificuldades técnicas e mudança de prioridades. Portanto, os resultados não demonstram que a implantação dessas capacidades de monitoramento tenha produzido os ganhos percebidos. Os próprios autores reconhecem a influência de outros fatores sobre o desempenho e a necessidade de reavaliar continuamente o perfil adotado \cite{plant2025}.

Linjama e Granlund~\citeano{linjama2026} investigaram, por estudo de caso qualitativo, a adoção de implantação contínua (\emph{continuous deployment}) em um projeto de cliente do setor público executado pela empresa finlandesa Solita. O projeto era iniciado do zero (\emph{greenfield}), sem sistemas legados, e o cliente já apresentava uma cultura DevOps estabelecida. Os participantes relataram liberação mais rápida de alterações, ciclos curtos de feedback, maior confiança nas atualizações e simplificação do diagnóstico pela utilização da mesma versão da aplicação em diferentes ambientes. Entre os desafios observados estavam instabilidade em testes de ponta a ponta, alterações de banco de dados em implantações sem interrupção e gerenciamento de \emph{feature flags}.

O estudo não identificou obstáculos à implantação contínua fundamentalmente exclusivos do setor público. Ressaltou, contudo, a relevância de requisitos de longa vida útil das aplicações, acessibilidade e gestão da informação. Por se tratar de um caso específico, seus resultados não podem ser generalizados automaticamente para instituições públicas com sistemas legados ou menor maturidade DevOps \cite{linjama2026}.

\subsubsection{Barreiras institucionais e tecnológicas da transformação digital governamental}

Figueredo Jr. e Silva~\citeano{figueredo2026}, pesquisadores vinculados à Universidade Federal da Grande Dourados, conduziram uma revisão sistemática das barreiras à implantação de serviços digitais na Administração Pública. O estudo utilizou ProKnow-C para selecionar o portfólio e o checklist/fluxograma PRISMA para documentar a revisão. A pesquisa partiu de 1.613 registros das bases Scopus e Web of Science e reuniu 85 estudos publicados entre 2019 e 2024.

As barreiras foram agrupadas em sete categorias: organizacionais/institucionais, técnicas/tecnológicas, humanas, socioculturais, econômicas, relativas a dados e segurança e legais. Os fatores organizacionais/institucionais apareceram em 68 dos 85 estudos (80\%), seguidos dos técnicos/tecnológicos em 59 e dos humanos em 58. O artigo também relaciona a predominância de determinados obstáculos ao nível de desenvolvimento dos países: restrições econômicas e de infraestrutura assumem maior relevância em contextos de IDH médio ou baixo, enquanto a integração de tecnologias avançadas apresenta outros desafios em contextos de maior desenvolvimento. Os autores advertem que sua estratégia de busca pode ter sub-representado a produção científica de países em desenvolvimento e menos desenvolvidos \cite{figueredo2026}.

No Brasil, o debate encontra referências normativas na Lei n.º 14.129/2021, que prevê princípios de Governo Digital, transparência e monitoramento da qualidade dos serviços públicos (art. 3.º, IV), além de painel de monitoramento do desempenho nas plataformas de Governo Digital (art. 20, II). A aplicação de seus comandos às administrações dos demais entes federados observa as condições do art. 2.º \cite{brasil2021}. O Decreto n.º 12.069/2024 estabelece, entre os objetivos da Estratégia Nacional de Governo Digital 2024--2027, ampliar a resiliência das estruturas tecnológicas e dispor de infraestrutura moderna, segura, escalável e robusta (art. 9.º, IV e VI). Quanto às Infraestruturas Públicas Digitais definidas pelo decreto, o art. 17, V, estabelece a priorização do mapeamento prévio de riscos e de medidas mitigadoras de privacidade, proteção de dados e segurança da informação \cite{brasil2024}.

Essas disposições contextualizam a importância institucional da qualidade e da governança dos serviços digitais. Não constituem, por si, comprovação empírica de que ferramentas de observabilidade tenham melhorado o desempenho de organizações públicas brasileiras.

\subsubsection{Observabilidade como capacidade de apoio à cultura DevOps}

Em mapeamento sistemático de 66 estudos publicados entre 2019 e 2024, Gomes, Rego e Trinta~\citeano{gomes2025}, pesquisadores da Universidade Federal do Ceará, investigaram fundamentos, classificações e desafios de observabilidade em aplicações baseadas em microsserviços. De acordo com o resumo da publicação, a modularidade dessas arquiteturas facilita aspectos do desenvolvimento, mas aumenta as dificuldades de operação e manutenção. Nesse contexto, os autores apontam limitações do monitoramento tradicional em aplicações nativas em nuvem e descrevem observabilidade como a capacidade de compreender e diagnosticar o comportamento interno de um sistema a partir de seus estados externos. O estudo propõe uma taxonomia e aborda a aplicação de observabilidade ao escalonamento automático de recursos \cite{gomes2025}.

Os fundamentos de coleta e análise de sinais operacionais discutidos na Seção~\ref{sec:observabilidade} podem ser relacionados a essa capacidade diagnóstica \cite{kosinska2023,giamattei2024}. Contudo, a passagem de informação técnica para colaboração organizacional exige análise cuidadosa: Plant et al.~\citeano{plant2025} identificaram monitoramento, documentação, gestão do conhecimento e feedback como capacidades pertinentes ao contexto de sua equipe pública, mas não avaliaram uma implantação completa de observabilidade. Gomes et al.~\citeano{gomes2025}, por sua vez, trataram de sistemas distribuídos, não da adoção de DevOps na Administração Pública.

A articulação entre essas pesquisas permite formular uma \emph{síntese teórica}: mecanismos de observabilidade podem fornecer informações úteis à investigação de falhas, ao acompanhamento da operação e à comunicação técnica entre equipes, oferecendo suporte potencial a práticas DevOps em organizações públicas. Isso não equivale a demonstrar causalmente redução de incidentes, eliminação de silos ou aumento de confiabilidade em órgãos públicos brasileiros \cite{plant2025,gomes2025,figueredo2026}.

A eventual adoção desses mecanismos também deve considerar riscos relativos ao tratamento de dados pessoais. Quando logs, rastreamentos ou outros dados operacionais permitirem identificar pessoas naturais, aplicam-se os princípios de finalidade, adequação, necessidade e segurança previstos no art. 6.º da Lei n.º 13.709/2018, bem como as regras de tratamento de dados pessoais pelo Poder Público previstas no art. 23 \cite{brasil2018}. Dessa forma, controles de acesso, limitação de coleta e gestão de retenção constituem medidas que devem ser avaliadas segundo a finalidade, os riscos e o contexto de cada sistema, não obrigações indistintas aplicáveis a toda telemetria.

Em síntese, a literatura oferece fundamentos para considerar a observabilidade uma capacidade técnica potencialmente útil ao fortalecimento de práticas DevOps na Administração Pública, desde que integrada ao contexto organizacional, à governança e às necessidades operacionais. Permanece, entretanto, uma lacuna na demonstração empírica dessa relação em organizações públicas brasileiras, a qual não pode ser suprida apenas pela nacionalidade dos pesquisadores ou pela transposição de estudos internacionais.


